% Chapter: Conclusions and Future Work

\section{Summary of Contributions}
\label{sec:conc_summary}

This thesis treated a structured Belousov--Zhabotinsky reaction--diffusion medium as the micro-scale picture inside a macro-scale porous chemical reactor.
The reactor's void space was abstracted as a pore-network graph: bulbous pores become nodes and narrow throats become edges.
The central claim was that the local pulse-traffic rules at each graph element can be measured in controlled two-dimensional simulations and then composed into a fast predictor of network-level behaviour.

Four local operators were isolated and quantified.
A straight throat is a near-perfect wire: the pulse speed is set by the kinetics and is almost independent of channel width.
Refractoriness limits the maximum signalling frequency; for the Oregonator working point used here, the absolute refractory period is 3.0 time units and the maximum one-to-one following rate is 0.167 per time unit.
A T-junction implements two-input inhibition when the readout is windowed in time: A AND (NOT B) was obtained with a separation ratio above $200\times$.
An anisotropic diffusion tensor steers pulses along preferred directions with a speed ratio that follows the eikonal $\sqrt{r}$ law.

These operators were encoded in an event-driven graph simulator.
The simulator correctly predicted the behaviour of an inhibition gate and a priority router, and it produced the correct truth table for a shared-control multi-junction gate.
Where it failed, the failure was instructive: a shared channel introduced crosstalk between two data paths that single-junction measurements did not capture, and a chambered Y-junction intended as a collision XOR leaked because the collision geometry allowed re-ignition of the output stem.
A thin-slab three-dimensional inhibition gate showed that the same operators transfer once the junction shape is corrected, and a one-way diode pilot showed that a simple asymmetric wall or static light barrier is not enough for robust rectification.
A light-field training attempt returned a negative result: within the optimisation budget, no improvement over a uniform suppression field was found.

The work therefore supports the following conclusion: network-level pulse routing in a structured excitable medium can be predicted from micro-scale transfer functions, but only when the geometry of every junction is carried explicitly in the model.
The graph abstraction is viable as a design layer between expensive continuum simulations and large network studies, but it must be calibrated against PDE solutions for each junction family.

\section{Limitations}
\label{sec:conc_limitations}

The results are limited to two-dimensional controlled simulations.
A real porous reactor is three-dimensional, and the interior light field is not known in the bulk, so the local kinetic parameters would be uncontrolled.
The 3D test reported here is a thin-slab boundary experiment, not a validated multi-gate network.

All measurements share the same computational parameter set, so the transfer functions are specific to that set.
The graph simulator uses scalar transit, refractory, and inhibition parameters extracted from idealised geometries; real junctions with rough walls, curved throats, or varying aspect ratios may require a larger library of local operators.

No experimental validation was attempted.
The mapping from model units to physical units is a calibration based on literature values, and the actual speeds and times in a wet experiment will depend on recipe, temperature, and illumination.

\section{Future Work}
\label{sec:conc_future}

The most immediate next step is to build a shape-resolved junction library.
Each entry would contain the measured transfer function for a family of junction geometries (T, Y, cross, chambered Y, angled branch), so the graph simulator can choose the right operator for each node rather than assuming a single universal value.

A second step is to incorporate uncertainty quantification into the graph simulator itself.
The one-at-a-time sensitivity analysis showed that kinetic parameters affect speed and amplitude; these uncertainties should propagate through the graph so that network-level predictions carry confidence intervals.

Third, the three-dimensional extension needs a volumetric illumination strategy.
Without a known $\phi(x,y,z)$ inside the bulk, the local operators cannot be assigned reliably.
Candidate approaches include structured light absorption, layered masks, or direct chemical patterning of the catalyst distribution.

Fourth, training should be reframed around geometry rather than around a small light-field grid.
The negative training result suggests that the native operators at the chosen working point are too stiff for small local modulation of the suppression field to change the routing decision.
Trainable porous processors may need trainable wall shape, catalyst patterning, or a richer parameterisation of the light field.

Finally, the framework should be tested against experimental data from a photosensitive BZ layer or a microfluidic excitable-medium network.
That would close the loop between the measured transfer functions and the pore-network graph, and would show whether the abstractions introduced here survive the noise and heterogeneity of a real chemical substrate.
